\documentclass[11pt]{article}

\usepackage[margin=1in]{geometry}
\usepackage{amsmath,amssymb,amsthm}
\usepackage{enumitem}
\usepackage[hidelinks]{hyperref}
\emergencystretch=5em

\newtheorem{theorem}{Theorem}
\newtheorem{lemma}{Lemma}

\newcommand{\N}{\mathbb N}
\newcommand{\Nzero}{\mathbb N_0}
\newcommand{\R}{\mathbb R}
\newcommand{\rhoF}{\mathcal P}
\newcommand{\omegaF}{\Omega}

\title{A Metric Counterexample: Nearly Isometric Embeddability Need Not Imply Almost Lipschitz Embeddability}
\author{Run \texttt{fa\_banach\_001}, agent \texttt{agent\_lane\_01}}
\date{2026-06-15}

\begin{document}
\maketitle

\section*{Source Problem}

This packet addresses Problem 4 in Section 5 of:

\begin{quote}
Florent Baudier and Gilles Lancien, \emph{Tight embeddability of proper and stable metric spaces}, arXiv:1503.06089.
\end{quote}

The source problem, on PDF page 17, asks:

\begin{quote}
Exhibit two metric spaces \(X\) and \(Y\) such that \(X\) nearly isometrically embeds into \(Y\), but \(X\) does not almost Lipschitz embed into \(Y\).
\end{quote}

\section*{Definitions Used}

We use the definitions from the source paper.  A metric space \(X\) almost
Lipschitz embeds into \(Y\) if there are constants \(r>0\), \(D\ge 1\), such
that for every continuous \(\varphi:[0,\infty)\to[0,1)\) with
\(\varphi(0)=0\) and \(\varphi(t)>0\) for \(t>0\), there is a map
\(f_\varphi:X\to Y\) satisfying
\[
 r\,d_X(x,y)\varphi(d_X(x,y))
 \le d_Y(f_\varphi(x),f_\varphi(y))
 \le D r\, d_X(x,y).
\]

For nearly isometric embeddability, the admissible compressions \(\rho\in\rhoF\)
are continuous functions with \(\rho(t)=t\) on \([0,1]\),
\(\rho(t)\le t\) for \(t\ge 1\), and \(\rho(t)/t\to0\).  The admissible
expansions \(\omega\in\omegaF\) satisfy \(\omega(0)=0\), \(t\le\omega(t)\) on
\([0,1]\), \(\omega(t)=t\) for \(t\ge1\), and \(\omega(t)/t\to\infty\) as
\(t\downarrow0\).  The space \(X\) nearly isometrically embeds into \(Y\) if
for every such pair \((\rho,\omega)\) there is \(f:X\to Y\) with
\[
\rho(d_X(x,y))\le d_Y(f(x),f(y))\le \omega(d_X(x,y)).
\]

\section*{Proof Intuition}

The source definition of nearly isometric embeddability permits the embedding
to depend on the requested sublinear compression \(\rho\).  We exploit this by
building a target \(Y\) with one ray for each admissible \(\rho\).  The
\(\rho\)-ray is metrically distorted by a sublinear function \(g_\rho\) that
still dominates \(\rho\), so the usual ray \(\Nzero\) can satisfy that
particular nearly isometric request by entering that component.

The obstruction to almost Lipschitz embeddability is global.  An almost
Lipschitz embedding of the integer ray would become a bi-Lipschitz embedding
after choosing a test function \(\varphi\) bounded below on all integer
distances.  But the bouquet target has no bi-Lipschitz ray.  A Lipschitz image
with linear lower growth must eventually be far from the common basepoint; at
that height it cannot switch components, because switching costs the sum of
the two heights.  Once trapped in one component, distances grow only
sublinearly, contradicting linear compression.

\section*{Auxiliary Concave Majorants}

\begin{lemma}
\label{lem:majorant}
Let \(\rho:[0,\infty)\to[0,\infty)\) be continuous, satisfy
\(\rho(t)=t\) on \([0,1]\), \(\rho(t)\le t\) for \(t\ge1\), and
\(\rho(t)/t\to0\) as \(t\to\infty\).  Then there is a continuous,
nondecreasing, concave, unbounded function \(g:[0,\infty)\to[0,\infty)\)
such that
\[
g(0)=0,\qquad g(t)=t\ (0\le t\le1),\qquad
\rho(t)\le g(t)\le t,\qquad \frac{g(t)}{t}\to0.
\]
Consequently \(d_g(m,n)=g(|m-n|)\) is a metric on \(\Nzero\).
\end{lemma}

\begin{proof}
Let \(h(t)=t\) for \(0\le t\le1\), and \(h(t)=1+\log t\) for \(t\ge1\).
Then \(h\) is continuous, nondecreasing, concave, unbounded,
\(h(t)\le t\), and \(h(t)/t\to0\).  Put \(\rho_0(t)=\max\{\rho(t),h(t)\}\).
Then \(\rho_0(t)\le t\), \(\rho_0(t)=t\) on \([0,1]\), and
\(\rho_0(t)/t\to0\).

Choose a sequence \(\varepsilon_j\downarrow0\).  Since
\(\rho_0(t)/t\to0\), for each \(j\) the number
\[
A_j=\sup_{t\ge0}\bigl(\rho_0(t)-\varepsilon_j t\bigr)
\]
is finite.  Define
\[
g(t)=\inf\Bigl(\{t\}\cup \{A_j+\varepsilon_j t:j\ge1\}\Bigr).
\]
Every affine function in the infimum dominates \(\rho_0\), and the function
\(t\mapsto t\) also dominates \(\rho_0\).  Hence \(g\ge\rho_0\ge\rho\).
Also \(g\le t\).  Because \(g\le t\) and \(g\ge\rho_0=t\) on \([0,1]\),
we have \(g(t)=t\) on \([0,1]\).

The pointwise infimum of increasing affine functions is concave and
nondecreasing.  For each fixed \(j\),
\[
\limsup_{t\to\infty}\frac{g(t)}{t}
\le \limsup_{t\to\infty}\frac{A_j+\varepsilon_j t}{t}
=\varepsilon_j,
\]
and therefore \(g(t)/t\to0\).  Finally, \(g\ge h\), so \(g\) is unbounded.

A nondecreasing concave function with \(g(0)=0\) is subadditive: concavity
implies \(g(s)/s\ge g(s+t)/(s+t)\) and \(g(t)/t\ge g(s+t)/(s+t)\) for
\(s,t>0\), whence \(g(s)+g(t)\ge g(s+t)\).  Thus \(g(|m-n|)\) satisfies the
triangle inequality on \(\Nzero\).
\end{proof}

\section*{The Counterexample}

\begin{theorem}
There are metric spaces \(X\) and \(Y\) such that \(X\) nearly isometrically
embeds into \(Y\), but \(X\) does not almost Lipschitz embed into \(Y\).
\end{theorem}

\begin{proof}
Let \(X=\Nzero\) with its usual metric \(d_X(m,n)=|m-n|\).
For every admissible \(\rho\in\rhoF\), choose a function \(g_\rho\) supplied by
Lemma~\ref{lem:majorant}.  Let \(Y_\rho=\Nzero\) with metric
\[
d_\rho(m,n)=g_\rho(|m-n|).
\]
Now form the metric wedge \(Y\) of the family \((Y_\rho)_{\rho\in\rhoF}\) by
identifying all zero points as a common basepoint \(o\).  If \(a\in Y_\rho\)
and \(b\in Y_\sigma\), set
\[
d_Y(a,b)=
\begin{cases}
g_\rho(|a-b|),& \rho=\sigma,\\
g_\rho(a)+g_\sigma(b),& \rho\ne\sigma.
\end{cases}
\]
This is the usual wedge metric: inside each component we use \(d_\rho\), and
between different components we travel through \(o\).

We first show that \(X\) nearly isometrically embeds into \(Y\).  Fix
\(\rho\in\rhoF\) and \(\omega\in\omegaF\).  Define \(F_{\rho,\omega}:X\to Y\)
by \(F_{\rho,\omega}(n)=n\) in the \(\rho\)-component.  If \(m,n\in\Nzero\)
and \(k=|m-n|\), then for \(k=0\) there is nothing to check.  For \(k\ge1\),
the defining properties of \(g_\rho\) and the fact that \(\omega(k)=k\) give
\[
\rho(k)\le g_\rho(k)
=d_Y(F_{\rho,\omega}(m),F_{\rho,\omega}(n))
\le k=\omega(k).
\]
Thus \(X\) nearly isometrically embeds into \(Y\).

It remains to show that \(X\) does not almost Lipschitz embed into \(Y\).
We first prove that \(Y\) contains no bi-Lipschitz copy of the usual ray.
Suppose, toward a contradiction, that \(F:\Nzero\to Y\) and constants
\(c,C>0\) satisfy
\[
c|m-n|\le d_Y(F(m),F(n))\le C|m-n|
\qquad(m,n\in\Nzero).
\]
Let \(R_n=d_Y(F(n),o)\).  Since
\[
R_n\ge d_Y(F(n),F(0))-R_0\ge cn-R_0,
\]
the heights \(R_n\) grow at least linearly.  On the other hand,
\(d_Y(F(n+1),F(n))\le C\).  If \(F(n)\) and \(F(n+1)\) lie in different
components, then
\[
d_Y(F(n+1),F(n))=R_{n+1}+R_n,
\]
which is eventually larger than \(C\).  Hence, for all sufficiently large
\(n\), consecutive points \(F(n)\) and \(F(n+1)\) lie in the same component.
The tail of the sequence \(F(n)\) is therefore contained in a single
component \(Y_{\rho_0}\).

Write \(F(n)=a_n\in Y_{\rho_0}\) for all \(n\ge N\).  Since
\[
g_{\rho_0}(|a_{n+1}-a_n|)=d_Y(F(n+1),F(n))\le C
\]
and \(g_{\rho_0}\) is nondecreasing and unbounded, there is \(B<\infty\) such
that \(|a_{n+1}-a_n|\le B\) for all \(n\ge N\).  Therefore
\[
|a_m-a_n|\le B|m-n| \qquad(m,n\ge N).
\]
For \(m,n\ge N\),
\[
d_Y(F(m),F(n))=g_{\rho_0}(|a_m-a_n|)
\le g_{\rho_0}(B|m-n|).
\]
Dividing by \(|m-n|\) and letting \(|m-n|\to\infty\), the right-hand side has
ratio tending to \(0\) because \(g_{\rho_0}(t)/t\to0\).  This contradicts the
lower bound \(d_Y(F(m),F(n))\ge c|m-n|\).  Thus no bi-Lipschitz embedding
\(\Nzero\to Y\) exists.

Finally suppose that \(X\) almost Lipschitz embedded into \(Y\).  Take
\[
\varphi(t)=
\begin{cases}
t/2,&0\le t\le1,\\
1/2,&t\ge1.
\end{cases}
\]
This \(\varphi\) is admissible in the almost Lipschitz definition.  Since all
nonzero distances in \(X=\Nzero\) are at least \(1\), the corresponding map
would satisfy
\[
\frac r2 |m-n|
\le d_Y(f_\varphi(m),f_\varphi(n))
\le Dr |m-n|
\qquad(m\ne n),
\]
which is a bi-Lipschitz embedding of the ray into \(Y\).  This contradicts the
previous paragraph.  Hence \(X\) does not almost Lipschitz embed into \(Y\).
\end{proof}

\section*{Verification Notes}

The construction is intentionally nonseparable in general, since it uses one
component for every admissible compression function \(\rho\).  The source
problem only asks for metric spaces, so no separability or Banach structure is
being claimed.  The proof uses no external theorem beyond elementary facts
about concave metric transforms and wedge metrics.

\section*{Bounded Novelty Check}

On 2026-06-15, the run indexes were searched for \texttt{1503.06089},
\texttt{nearly isometric}, and \texttt{almost Lipschitz}; no duplicate packet
or attempt for this metric problem was found.  Local parsed-source search for
phrases such as \texttt{nearly isometrically embeds} and
\texttt{does not almost Lipschitz} found the source paper but no separate
answer.  Web searches for exact phrases including
\texttt{"nearly isometrically embeds" "almost Lipschitz"},
\texttt{"Exhibit two metric spaces" "nearly isometrically" "almost Lipschitz"},
and \texttt{site:arxiv.org "almost Lipschitz" "nearly isometric"} returned
the source paper and no later answer.  The novelty claim is therefore
bounded by these searches.

\section*{References}

\begin{itemize}[leftmargin=*]
\item F. Baudier and G. Lancien, \emph{Tight embeddability of proper and stable metric spaces}, arXiv:1503.06089.
\end{itemize}

\end{document}
